% !TeX program = pdflatex
% =============================================================================
% Zusatzaufgaben: Gewichtskraft & Federkraft
% UZH Praktikum I -- Sara Romer Urban, Kantonsschule im Lee, HS2026
% Klasse: 2e (26.08.2026)
%
% Reine Uebungsaufgaben fuer die Vor-/Nachbereitung (nicht fuer den
% Lektioneneinsatz selbst) -- siehe arbeitsblatt2-gewichtskraft-federkraft.tex
% (Theorie + Plenum) und aufgabenblatt1-gewichtskraft-federkraft.tex
% (Stationsexperiment, Partnerarbeit) fuer den Lektioneneinsatz.
% Aufgabe 1 (Gewichtskraft berechnen) und Aufgabe 2 (Einheiten der
% Federkonstante) waren urspruenglich im Arbeitsblatt, wurden von dort
% hierher verschoben. Aufgaben 3-5 sind neu, adaptiert aus den
% LEIFIphysik-Referenzquellen in diesem Ordner (Astronautenwissen,
% Gravitation auf unbekannten Planeten, Feder auf dem Mond) -- per
% lernziele.md als "Ergaenzend als Uebungsaufgaben" vorgesehen.
% Aufgabe 5 nutzt bereits F=D*Delta x (Hookesches Gesetz): unproblematisch,
% da Zusatzaufgaben erst nach der Lernaufgabe im Aufgabenblatt bearbeitet
% werden, wenn die SuS die Formel schon selbst hergeleitet haben.
% =============================================================================
\documentclass[addpoints,11pt,a4paper]{exam}

\usepackage[T1]{fontenc}
\usepackage[utf8]{inputenc}
\usepackage[ngerman]{babel}
\usepackage{lmodern}
\usepackage[a4paper,margin=2.2cm,headheight=14pt]{geometry}
\usepackage{amsmath}
\usepackage{siunitx}
\sisetup{per-mode=fraction,fraction-command=\frac}
\usepackage{array,booktabs,tabularx,multirow}
\usepackage{enumitem}
\usepackage{xcolor}
\usepackage{colortbl}
\usepackage{tikz}
\usepackage{physics}
\usepackage[most]{tcolorbox}
\usepackage{lastpage}

\usetikzlibrary{arrows.meta,calc,angles,quotes}

\usepackage{setspace}
\setlength{\parindent}{0pt}
\setlength{\parskip}{0.6em}
\setstretch{1.08}
\setlist[itemize]{itemsep=0.4em}

% -----------------------------
% Farben (Corporate-Farbschema Physik KFR)
% -----------------------------
\definecolor{titleblue}{HTML}{183A6B}
\definecolor{accentblue}{HTML}{2E6F9E}
\definecolor{lightblue}{HTML}{EAF4FB}
\definecolor{verylightblue}{HTML}{F7FBFE}
\definecolor{lightgray}{HTML}{F4F6F8}
\definecolor{darkgray}{HTML}{4A4A4A}
\definecolor{accentpurple}{HTML}{7B2D8E}
\definecolor{groupteal}{HTML}{1B7F72}
\definecolor{lightgroupteal}{HTML}{E8F5F3}
\definecolor{lightpurple}{HTML}{F3EAF7}

\newcolumntype{Y}{>{\centering\arraybackslash}X}
\newcolumntype{P}[1]{>{\raggedright\arraybackslash}p{#1}}

\tcbset{before skip=10pt, after skip=10pt}
\newtcolorbox{titlebox}{
  enhanced, colback=titleblue, colframe=titleblue, boxrule=0pt, arc=3mm,
  left=3mm,right=3mm,top=3mm,bottom=3mm
}

\newlength{\aufgabenindent}
\settowidth{\aufgabenindent}{10.\hskip\labelsep}

% Praktikum I: Lernziele/Einleitung/Theorie/Aufgaben werden NICHT mehr
% farbig gerahmt (nur Formel-, Hinweis- und Antwortbox bleiben Boxen) --
% siehe institutionen/UZH-Praktikum-I/CLAUDE.md, Abschnitt "Boxen-Layout".
\newenvironment{infobox}{%
  \par\addvspace{10pt}\noindent
}{%
  \par\addvspace{10pt}
}

% taskbox/groupbox stehen in questions VOR dem ersten \question (Bilder,
% Fliesstext) -- exam.cls' questions-Liste braucht dafür eine echte,
% eigenständige Box (sonst "Something's wrong--perhaps a missing \item",
% weil z.B. ein \begin{center} vor dem ersten \item der Aussenliste steht).
% Deshalb bleiben es tcolorboxen, nur unsichtbar gemacht (keine Farbe/Rahmen,
% kein Padding) statt sie durch reine Absätze zu ersetzen.
\newtcolorbox{taskbox}[1][]{
  enhanced, breakable,
  colback=white, colframe=white, boxrule=0pt,
  left=0mm,right=0mm,top=0mm,bottom=0mm,
  before skip=10pt, after skip=10pt,
  before upper={%
    \def\tmpboxtitle{#1}%
    \ifx\tmpboxtitle\empty\else\noindent\textbf{#1}\par\smallskip\fi
  }
}

\newtcolorbox{groupbox}[1][Gruppenarbeit]{
  enhanced, breakable,
  colback=white, colframe=white, boxrule=0pt,
  left=0mm,right=0mm,top=0mm,bottom=0mm,
  before skip=10pt, after skip=10pt,
  before upper={\noindent\textbf{#1}\par\smallskip}
}

% formelbox: kurze Formelsammlung zum Nachschlagen -- kein Theorie-Ersatz,
% keine Herleitung (siehe arbeitsblatt2-gewichtskraft-federkraft.tex
% dafuer), nur die fertigen Formeln als Gedaechtnisstuetze.
\newtcolorbox{formelbox}[1][Formel]{
  enhanced, breakable, colback=lightpurple, colframe=accentpurple,
  title={#1}, fonttitle=\bfseries, boxrule=0.9pt, arc=2mm,
  left=5mm,right=5mm,top=2mm,bottom=2mm
}

\newtcolorbox{answerbox}{
  breakable, colback=white, colframe=gray!55, boxrule=0.5pt, arc=1.5mm,
  left=2mm,right=2mm,top=2mm,bottom=2mm
}

% --- Lösungen ein-/ausblenden -----------------------------------------------
\noprintanswers
% \printanswers

\pointformat{}
\renewcommand{\solutiontitle}{\noindent\color{titleblue}\textbf{Lösung:}\enspace}

\pagestyle{headandfoot}
\cfoot{\textcolor{darkgray}{Seite \thepage\ von \pageref{LastPage}}}

\newcommand{\Kopfzeile}[4]{%
  \begin{titlebox}
    {\color{white}\sffamily\bfseries\Large #4}
  \end{titlebox}
  \vspace{0.5em}
  \noindent\textbf{Klasse:}~#1 \hfill \textbf{Datum:}~#2 \hfill \textbf{Thema:}~#3
    \hfill \textbf{Lehrperson:}~Loris Keller
  \vspace{0.8em}
  \lhead{\textcolor{darkgray}{#3}}
  \rhead{\textcolor{darkgray}{#4}}
}

\newlength{\antwortraumlen}
\newcommand{\Antwortraum}[1]{%
  \pgfmathsetlength{\antwortraumlen}{#1*2.5cm}%
  \begin{answerbox}
    \rule{0pt}{\antwortraumlen}%
  \end{answerbox}%
}

\newcommand{\Klasse}{2e}
\newcommand{\Datum}{26.08.2026}
\newcommand{\Thema}{Gewichtskraft \& Federkraft}
\newcommand{\ZusatzaufgabenTitel}{Zusatzaufgaben: Gewichtskraft \& Federkraft}

\begin{document}

\Kopfzeile{\Klasse}{\Datum}{\Thema}{\ZusatzaufgabenTitel}

\begin{infobox}
\textbf{Zu diesen Zusatzaufgaben}\\
Diese Übungsaufgaben ergänzen das Arbeitsblatt und das Aufgabenblatt
\emph{Gewichtskraft \& Federkraft} und sind für die Vor-/Nachbereitung
gedacht, nicht für den Lektioneneinsatz selbst. Alle benötigten Formeln
und Herleitungen finden Sie im Arbeitsblatt bzw. im Aufgabenblatt.
\end{infobox}

\begin{formelbox}[Formelsammlung]
\[
  F_G = m\cdot g, \qquad F = D\cdot\Delta x.
\]
$g$ in $\si{\meter\per\second\squared}$ oder gleichwertig in
$\si{\newton\per\kilogram}$ (es gilt $1\,\si{\newton\per\kilogram} =
1\,\si{\meter\per\second\squared}$); $F_G$, $F$ in $\si{\newton}$; $m$ in
$\si{\kilogram}$; $D$ in $\si{\newton\per\meter}$ (oder
$\si{\newton\per\centi\meter}$); $\Delta x$ in $\si{\meter}$ oder
$\si{\centi\meter}$.
\end{formelbox}

\begin{questions}

\begin{taskbox}[Aufgabe 1: Gewichtskraft berechnen]
\question Eine Astronautin hat eine Masse von $m=\SI{72}{\kilogram}$
(inklusive Ausrüstung).
\begin{parts}
  \part[2] Berechnen Sie ihre Gewichtskraft auf der Erde.
  \Antwortraum{2}
  \begin{solution}
  \[
    F_{G,\text{Erde}} = m\cdot g_{\text{Erde}} = \SI{72}{\kilogram}\cdot\SI{9.81}{\meter\per\second\squared} = \SI{706.3}{\newton}.
  \]
  \end{solution}

  \part[2] Berechnen Sie ihre Gewichtskraft auf dem Mond und vergleichen Sie
  das Ergebnis mit (a). Was ist an ihrer Masse in diesem Vergleich anders,
  was bleibt gleich?
  \Antwortraum{2}
  \begin{solution}
  \[
    F_{G,\text{Mond}} = m\cdot g_{\text{Mond}} = \SI{72}{\kilogram}\cdot\SI{1.62}{\meter\per\second\squared} = \SI{116.6}{\newton}.
  \]
  Die Gewichtskraft ist auf dem Mond rund sechsmal kleiner als auf der
  Erde. Ihre Masse $m=\SI{72}{\kilogram}$ bleibt dabei in beiden Fällen
  exakt gleich. Nur der Ortsfaktor $g$ hat sich geändert.
  \end{solution}

  \part[2] Eine Marssonde mit der Masse $m=\SI{150}{\kilogram}$ wird nach
  der Landung auf dem Mars gewogen und zeigt eine Gewichtskraft von
  $F_G=\SI{556.5}{\newton}$. Bestimmen Sie daraus den Ortsfaktor $g_{\text{Mars}}$.
  \Antwortraum{2}
  \begin{solution}
  Aus $F_G=m\cdot g$ nach $g$ aufgelöst:
  \[
    g_{\text{Mars}} = \frac{F_G}{m} = \frac{\SI{556.5}{\newton}}{\SI{150}{\kilogram}} = \SI{3.71}{\meter\per\second\squared}.
  \]
  Der Ortsfaktor auf dem Mars liegt damit zwischen dem des Mondes und dem
  der Erde. Er entspricht dem tatsächlichen Wert für den Planeten Mars.
  \end{solution}
\end{parts}
\end{taskbox}

\begin{taskbox}[Aufgabe 2: Einheiten der Federkonstante]
\question Zwei Federn wurden vermessen. Aus den jeweiligen
$F$-$\Delta x$-Diagrammen wurden folgende Steigungen abgelesen:
Feder~1: $D_1=\SI{42}{\newton\per\centi\meter}$,
Feder~2: $D_2=\SI{25}{\newton\per\centi\meter}$.
\begin{parts}
  \part[2] Rechnen Sie beide Werte in die SI-Einheit N/m um. Zeigen Sie Ihren
  Rechenweg.
  \Antwortraum{2}
  \begin{solution}
  Da $\SI{1}{\centi\meter}=\SI{0.01}{\meter}$, gilt $\SI{1}{\newton\per\centi\meter}=\SI{100}{\newton\per\meter}$:
  \[
    D_1 = \SI{42}{\newton\per\centi\meter} = 42\cdot\SI{100}{\newton\per\meter} = \SI{4200}{\newton\per\meter}.
  \]
  \[
    D_2 = \SI{25}{\newton\per\centi\meter} = 25\cdot\SI{100}{\newton\per\meter} = \SI{2500}{\newton\per\meter}.
  \]
  \end{solution}

  \part[2] Welche der beiden Federn ist die stärkere (härtere) Feder, und
  wie erkennt man das direkt am Wert von $D$? Ändert sich diese Antwort,
  wenn man statt N/m die Werte in N/cm vergleicht?
  \Antwortraum{2}
  \begin{solution}
  Feder~1 ist die stärkere Feder, da sie den grösseren Wert von $D$ besitzt:
  Eine grössere Federkonstante $D$ bedeutet, dass bei gleicher Kraft eine
  kleinere Längenänderung $\Delta x$ auftritt, die Feder also \glqq härter\grqq{}
  ist. Der Vergleich \glqq Feder~1 ist stärker als Feder~2\grqq{} ändert sich
  nicht, egal ob man N/cm- oder N/m-Werte vergleicht: Die Einheit
  beeinflusst nur die Zahlenwerte, nicht die Rangfolge der beiden Federn.
  \end{solution}
\end{parts}
\end{taskbox}

\begin{taskbox}[Aufgabe 3: Astronautenwissen]
\question Ein Astronaut befindet sich weit entfernt von jedem
Himmelskörper im Weltall, vollständig schwerelos.
\begin{parts}
  \part[2] Kann er dort mit einer Federwaage oder mit einer Balkenwaage das
  Gewicht eines Gegenstands bestimmen? Begründen Sie Ihre Antwort.
  \Antwortraum{2}
  \begin{solution}
  Nein, mit keiner der beiden Waagen: Der Gegenstand ist schwerelos, es
  gibt also gar keine Gewichtskraft, die eine der beiden Waagen anzeigen
  könnte. Die Masse liesse sich allenfalls über einen
  Beschleunigungsversuch bestimmen (vgl. die Übung \glqq Wiegen im
  Weltraum\grqq{} im Arbeitsblatt).
  \end{solution}

  \part[4] Der Astronaut hat vor dem Start zum Mond eine Masse von
  $m=\SI{72.0}{\kilogram}$. Durch die Anstrengung des Fluges nimmt er ab.
  Mit einer Federwaage stellt er auf dem Mond ein Körpergewicht von
  $F_G=\SI{115}{\newton}$ fest. Berechnen Sie, wie viele Kilogramm er
  abgenommen hat. Rechnen Sie mit $g_{\text{Erde}}=\SI{9.81}{\newton\per\kilogram}$
  und $g_{\text{Mond}}=\SI{1.62}{\newton\per\kilogram}$.
  \Antwortraum{4}
  \begin{solution}
  Seine neue Masse auf dem Mond folgt aus $F_G=m\cdot g_{\text{Mond}}$,
  aufgelöst nach $m$:
  \[
    m = \frac{F_G}{g_{\text{Mond}}} = \frac{\SI{115}{\newton}}{\SI{1.62}{\newton\per\kilogram}} = \SI{71.0}{\kilogram}.
  \]
  Abgenommen hat er also
  \[
    \SI{72.0}{\kilogram} - \SI{71.0}{\kilogram} = \SI{1.0}{\kilogram}.
  \]
  \end{solution}
\end{parts}
\end{taskbox}

\begin{taskbox}[Aufgabe 4: Gravitation auf einem unbekannten Planeten]
\question Ein Astronaut wiegt auf der Erde für die Reise zu einem
unbekannten Planeten einen Reisvorrat von $m=\SI{20}{\kilogram}$ ab.
Rechnen Sie in dieser Aufgabe mit $g_{\text{Erde}}=\SI{10}{\newton\per\kilogram}$.
\begin{parts}
  \part[2] Der Astronaut landet mit seinem Vorrat auf dem unbekannten
  Planeten. Was kann er ohne zusätzliche Hilfsmittel über Masse und
  Gewichtskraft seines Vorrats auf dem Planeten aussagen?
  \Antwortraum{2}
  \begin{solution}
  Die Masse ist ortsunabhängig und beträgt daher weiterhin
  $\SI{20}{\kilogram}$. Über die Gewichtskraft kann er ohne weitere
  Hilfsmittel dagegen keine Aussage machen, da der Ortsfaktor auf dem
  unbekannten Planeten nicht bekannt ist.
  \end{solution}

  \part[2] Welches Hilfsmittel bräuchte der Astronaut, um mithilfe seines
  Vorrats den Ortsfaktor auf dem unbekannten Planeten zu bestimmen?
  \Antwortraum{2}
  \begin{solution}
  Er bräuchte eine Federwaage, mit der er die Gewichtskraft des Vorrats auf
  dem Planeten messen könnte. Zusammen mit der bekannten Masse liesse sich
  daraus über $g=F_G/m$ der Ortsfaktor bestimmen.
  \end{solution}

  \part[4] Der Astronaut verzehrt ein Drittel seines Vorrats. Zufällig ist
  die Gewichtskraft des restlichen Vorrats auf dem unbekannten Planeten
  jetzt genau so gross wie die Gewichtskraft des vollen Vorrats auf der
  Erde. Bestimmen Sie den Ortsfaktor $g_x$ auf dem unbekannten Planeten.
  \Antwortraum{4}
  \begin{solution}
  Die Gewichtskraft des vollen Vorrats auf der Erde:
  \[
    F_{G,\text{Erde}} = m\cdot g_{\text{Erde}} = \SI{20}{\kilogram}\cdot\SI{10}{\newton\per\kilogram} = \SI{200}{\newton}.
  \]
  Die Restmasse nach dem Verzehr eines Drittels beträgt
  \[
    m^\ast = \frac{2}{3}\cdot m = \frac{2}{3}\cdot\SI{20}{\kilogram} \approx \SI{13}{\kilogram}.
  \]
  Da die Gewichtskraft des Restvorrats auf dem unbekannten Planeten
  ebenfalls $\SI{200}{\newton}$ beträgt, folgt
  \[
    g_x = \frac{F_{G,x}}{m^\ast} = \frac{\SI{200}{\newton}}{\SI{13}{\kilogram}} \approx \SI{15}{\newton\per\kilogram}.
  \]
  \end{solution}
\end{parts}
\end{taskbox}

\begin{taskbox}[Aufgabe 5: Feder auf dem Mond]
\question Eine elastische Schraubenfeder mit der Federkonstante
$D=\SI{3.00}{\newton\per\centi\meter}$ wird auf dem Mond mit einem Körper
der Masse $m=\SI{5.00}{\kilogram}$ belastet. Rechnen Sie mit
$g_{\text{Mond}}=\SI{1.62}{\newton\per\kilogram}$.
\begin{parts}
  \part[2] Berechnen Sie die Gewichtskraft des Körpers auf dem Mond.
  \Antwortraum{2}
  \begin{solution}
  \[
    F_G = m\cdot g_{\text{Mond}} = \SI{5.00}{\kilogram}\cdot\SI{1.62}{\newton\per\kilogram} = \SI{8.10}{\newton}.
  \]
  \end{solution}

  \part[2] Berechnen Sie, um wie viel sich die Feder dabei verlängert.
  \Antwortraum{2}
  \begin{solution}
  Im Gleichgewicht gilt $F=F_G$, also $D\cdot\Delta x=F_G$, aufgelöst nach
  $\Delta x$:
  \[
    \Delta x = \frac{F_G}{D} = \frac{\SI{8.10}{\newton}}{\SI{3.00}{\newton\per\centi\meter}} = \SI{2.70}{\centi\meter}.
  \]
  \end{solution}
\end{parts}
\end{taskbox}

\end{questions}

\end{document}
